\section{Conclusion}


To fix the problem of 3D data scarcity in zero-shot 3D classification, we propose TeGA (Text-guided Geometric Augmentation), a language-image-3D dataset expansion method with high-quality synthetic 3D data. TeGA plays a significant role in the expansion of existing datasets.
To further enhance the quality of 3D data, we apply consistency filtering to remove unaligned data that could negatively impact training. In practice, by combining the generated synthetic data with existing 3D datasets, we demonstrated improved performance in zero-shot 3D classification. 
Based on these results, we believe text-to-3D models have the potential to alleviate challenges associated with 3D data collection and, when leveraged for data expansion, can support the development of generalizable vision models for 3D vision.

\noindent{\textbf{Limitations.}} In this paper, we adopted Point-E as a text-to-3D model. However, Point-E is a generative model trained on data that is not open to the public. Since it is difficult to generate out-of-distribution data not included in the training dataset, our framework heavily relies on the generative model. Therefore, we plan to focus on improving the text-to-3D model in the future.
Additionally, when incorporating synthetic 3D data generated by Point-E into the training dataset, there is concern that societal biases or other unintended elements may be inadvertently introduced. It is recommended to take these risks into account and incorporate the findings from this study into the model improvement process. Furthermore, we look forward to the development of clean Text-to-3D models, like Stable Diffusion, which are built on publicly available datasets.
